\section{规模化训练首先是资源流问题}

“一万张 GPU”不是把单卡程序复制一万份。模型参数、梯度、优化器状态、激活、输入数据和 checkpoint 会在不同存储层与设备之间移动；任何一条流量超过链路能力，整次训练就会在等待中停住。本课的主线因此不是记住五个 parallelism 名称，而是对每一种方案反复追问：\textbf{切分什么、复制什么、用哪个 collective 恢复等价计算、通信能否被计算隐藏？}

\subsection{为什么大模型仍在继续变大}

开场三页把“规模”拆成模型、数据和上下文三个轴。第一页是近期模型地图，第二页把参数量与能力趋势并列，第三页给出课堂快照中的数量级：trillion-level parameters、约 $15$T training tokens 和 million-token context。它们不是统一配方，却说明单一 GPU 的容量、带宽和迭代时间会同时受到压力。

\lecturefigure{slide-004.jpg}{模型发布趋势：训练规模仍在扩大}{CS25 V6 official deck, p.~4}

\lecturefigure{slide-005.jpg}{模型规模与训练计算的公开趋势对照}{CS25 V6 official deck, p.~5}

\lecturefigure{slide-006.jpg}{课堂数量级快照：参数、训练 token 与上下文长度}{CS25 V6 official deck, p.~6}

\begin{knowledgebox}{读图：三个规模轴不能互相替代}
参数量决定持久模型状态和每层矩阵计算的基线，训练 token 数决定总迭代预算与输入管线压力，上下文长度决定单样本 activation、attention working set 和可用 micro-batch。增加 DP 只能更快消费更多 batch，不能让单个过长 sequence 自动装下；增加 CP 能切长 sequence，却不能减少完整模型状态；ZeRO（Zero Redundancy Optimizer，零冗余优化）能把 optimizer state、gradient、parameter 按 stage 分片，却不改变总训练 token。读后续每种 parallelism 时，都应先把它放回这三个轴，判断它究竟解除哪个约束，又把成本转移到了哪里。
\end{knowledgebox}

\teachervoice{00:00:48--00:03:02，Tazi 说明本课以 Ultra-Scale Playbook 为基础，但补入更新的 MoE scaling 实践。他把压力同时归因于数据读取、每步迭代、显存和 checkpoint 写入，而不是只谈参数量。}

\begin{warningbox}{规模与智能相关，不等于规模单独造成智能}
课堂用近期模型说明“训练更大”仍是产业趋势，但相关性不能隔离数据质量、架构、后训练、工具使用和评测污染。系统讲义真正需要继承的是资源数量级与工程约束，而不是把某个参数规模写成能力定律。
\end{warningbox}

\subsection{一次训练迭代到底搬运什么}

训练基础设施至少要持续读取 token 流，把模型与状态放进加速器内存，在 forward/backward 之间保留激活，执行 optimizer update，并定期把可恢复状态写回存储。下图中的 ``$\sim15$T tokens read''、``$\sim1$s iteration''、``$\sim80$GB per GPU'' 和 ``$\sim1$TB checkpoint'' 是课堂示意量级；读图时应把它们理解为四类独立瓶颈，而不是同一模型的精确配置。

\lecturefigure{slide-009.jpg}{LLM 训练对数据、计算、显存与 checkpoint I/O 的联合压力}{CS25 V6 official deck, p.~9}

\begin{knowledgebox}{读图：把一次 step 看成多级流水线}
训练 step 开始前，data loader 要从存储和主机内存准备样本；GPU 随后读取参数与 activation，执行 kernel，并通过互连同步；step 末尾 optimizer 写回参数，而 checkpoint 周期还会把更大的恢复状态送到持久存储。任何一段吞吐不足都会产生不同形状的 stall：输入不足表现为 GPU 开始前空闲，collective 过慢表现为 backward/forward 间等待，checkpoint 过慢表现为周期性长暂停。定位问题时应先确定 stall 属于哪一级，而不是统一归因于“GPU 不够快”。
\end{knowledgebox}

\begin{knowledgebox}{术语消化：第一组资源对象}
\textbf{HBM（High Bandwidth Memory，高带宽内存）}是 GPU 附近的高带宽 DRAM，常被口语化称为显存；它容量有限但供给 kernel 的带宽很高。\textbf{Optimizer state（优化器状态）}是 Adam/AdamW 等算法为每个参数维护的一阶矩 $m$、二阶矩 $v$ 等更新统计，常比参数本身占用更多字节。\textbf{Activation（激活）}是 forward 中间结果，backward 需要它来计算梯度。三者必须分开记账。
\end{knowledgebox}

\subsection{单卡 baseline 与 gradient accumulation 的隐藏假设}

单卡训练的逻辑顺序很清楚：取 batch，做 forward，计算 loss，backward 得到 gradients，optimizer 读取参数与状态并写回更新后的模型。当 global batch size（GBS，全局批大小）远大于单卡可容纳的 micro-batch 时，可以多次 forward/backward 累加梯度，再执行一次更新。

\lecturefigure{slide-010.jpg}{单 GPU 的 forward--backward--optimizer 基线}{CS25 V6 official deck, p.~10}

\lecturefigure{slide-011.jpg}{用 gradient accumulation 构造更大的 global batch}{CS25 V6 official deck, p.~11}

但 gradient accumulation 只解决“一个更新由多少样本组成”，并不自动解决模型是否装得下，也不会把串行 micro-batch 变快。下页把这两个前提明确写出：训练状态必须先能驻留在单卡；如果等待所有 micro-batch 顺序完成，设备数量再多也没有被利用。

\lecturefigure{slide-012.jpg}{单卡累积的两个前提：状态能装下，而且串行等待可接受}{CS25 V6 official deck, p.~12}

设一次更新包含 $M$ 个 micro-batch，每个 micro-batch 的 loss 为 $\ell_j(\theta)$，则常见的累积梯度写成
\[
g=\frac{1}{M}\sum_{j=1}^{M}\nabla_\theta \ell_j(\theta),
\qquad
\theta' = \operatorname{Opt}(\theta,g,s).
\]
其中，$M$ 是累积步数，$\theta$ 是参数，$g$ 是本次更新使用的平均梯度，$s$ 是 optimizer state，$\operatorname{Opt}$ 是更新规则，$\theta'$ 是新参数。这个公式没有规定 $M$ 个梯度必须在同一设备上产生；分布式训练正是重新安排这些计算与状态的所有权。

\subsection{五个切分轴与统一判断框架}

课程把方案分为 Data、Tensor/Sequence、Pipeline、Context 和 Expert 五类。它们分别沿 batch、hidden dimension、layer、sequence 和 expert 轴切分。开场的 agenda 与 preamble 还强调：这些机制可用于 GPU 或 TPU，也不局限于 pretraining；从两张设备开始，资源切分就可能有意义。

\lecturefigure{slide-013.jpg}{课程路线：从 DP/ZeRO 到 TP、PP、CP 与 EP}{CS25 V6 official deck, p.~13}

这里首次出现的 \term{ZeRO（Zero Redundancy Optimizer，零冗余优化）}不是一种新 optimizer，而是把 data-parallel ranks 上原本重复的 optimizer state、gradient 和 parameter 按 stage 逐步分片的状态管理方法。ZeRO-1/2/3 分别增加被分片的状态种类；Adam/AdamW 的一阶矩 $m$ 与二阶矩 $v$ 仍按原公式更新，只是所有权和收集时机改变。

\lecturefigure{slide-014.jpg}{适用边界：多种加速器、训练阶段与设备规模}{CS25 V6 official deck, p.~14}

所谓 \term{sharding（分片）}，是把一个原本在每张设备上完整复制的对象切成若干 shard，并让不同 rank 只拥有其中一部分。必须同时说明“切的是什么”和“沿哪个设备组切”：参数分片、梯度分片、batch 分片、sequence 分片并不是一回事。

\lecturefigure{slide-015.jpg}{目标一：分片模型、梯度与 optimizer state}{CS25 V6 official deck, p.~15}

\lecturefigure{slide-016.jpg}{目标二：分片数据并并行处理 micro-batch}{CS25 V6 official deck, p.~16}

\begin{importantbox}{统一四问}
分析任何 parallelism 时依次回答：一，哪个张量维度或训练状态被切分；二，哪些对象仍被复制；三，哪个集合通信让结果与单设备数学等价；四，该通信位于 critical path 还是能与 compute overlap。若这四问说不清，所谓“显存节省”通常只是把成本转移到了网络或调度器。
\end{importantbox}

\begin{knowledgebox}{背景概念：通信成本不只有字节数}
一次 collective 的时间可粗略写成 $T\approx \alpha n_{\mathrm{steps}}+\beta V$：$\alpha$ 表示每轮启动与同步延迟，$n_{\mathrm{steps}}$ 取决于 ring/tree 等算法和 rank 数，$\beta$ 是每字节传输时间，$V$ 是实际流量。小消息常受 $\alpha$ 主导，大消息更受带宽和拓扑拥塞主导；跨节点消息还会经过不同 NIC、交换机与路由。于是同样的总字节，拆成很多细碎 all-gather 可能比少量大消息更慢；反过来，大消息若启动太晚又无法 overlap。并行设计需要同时优化消息大小、次数、参与组和启动时机。
\end{knowledgebox}

\subsection{本章小结}

规模化训练的输入不是“GPU 数量”，而是一组资源约束：状态容量、激活容量、global batch、sequence length、模型结构、链路拓扑和存储吞吐。五类并行分别重排这些资源的所有权；后续章节将从最直观的 data parallelism 开始，逐步增加通信与调度复杂度。

